% ============================================================
% 7. DISCUSSION
% ============================================================
\section{Discussion}

% [REVISION | §7.1-Implications-trim-Day10-batch3 | 2026-05-15 | Day-10 Batch 3 page-budget cut #11: trimmed §7.1 from 4 paragraphs (opening framing + Monitor-design + Reward-design + Unifying-design-lesson) to 2 paragraphs (Reward-design specifics + Generalized design lesson). The opening framing paragraph and the Monitor-design paragraph are dropped: under Batch 5 drop-C2 the monitor-design framing loses its empirical anchor (we make no monitor-driven Dimensional Escape claim in body), and the dropped paragraphs duplicate the Definition~\ref{def:dim_escape} synthesis. The Unifying-design-lesson paragraph is preserved (50\% shorter) as the generalization vehicle. The 2026-05-11 §7.1-Implications-recast-D+ marker is subsumed; the white-box / interpretability-tool framing carried by the Reward-design paragraph is preserved. Per Day-9 page-budget cut plan §3 rank #11 (MEDIUM risk; touches Contribution 3 framing, save 0.30 pp). — pending audit]
\subsection{Implications for Latent-Space Reward Design}\label{sec:discussion-implications}

The reward-driven manifestation of Dimensional Escape---Hypothesis~\ref{hyp:goodhart}, empirically established in §\ref{sec:exp1}---carries a direct implication for the design of test-time scaling procedures that operate on continuous latent representations. A reward signal validated as a static correlate of correctness ($|r| \approx 0.60$ between next-token entropy and ARC-Easy answer correctness) does not survive deployment as an active MCTS optimization objective: at 1B the search amplifies the model's first-letter prior; at 3B it collapses to greedy decoding. Neither pathway corresponds to ``find the latent state that produces the correct answer''---both are degenerate optima that satisfy the entropy proxy without traversing the semantically relevant subspace. Anti-Goodhart reward designs surveyed in §\ref{sec:future-work}---multi-channel ensembles, external Process Reward Models, self-consistency across stochastic rollouts, counterfactual sensitivity, anti-prior penalties, information-bottleneck objectives---all attempt to raise the effective rank of the scoring projection or to anchor it in exogenous signal that the search cannot itself rewrite. Whether any resists Dimensional Escape empirically is an open question; our Experiment~1 result establishes the null-control against which any candidate must be measured.

The design response is to raise the effective rank of the scoring projection (e.g., attention-aggregated probes) or to ground the scoring in signal the searched representation cannot rewrite (exogenous answer verifiers, ground-truth-matched reward functions); Reversible KV-Cache MCTS provides white-box infrastructure for evaluating which response class generalizes.

\subsection{Broader Impact}

% [REVISION | §7.2-BroaderImpact-recast-D+ | 2026-05-11 | Wholesale rewrite of §7.2 Broader Impact dropping α+/Candidate-A "discover alignment-breaking interventions" framing and the three associated "mitigations" (controlled offline environment, cryptographic-hash audit trail, responsible-disclosure of vulnerabilities) — all premised on the false claim that the work produces deployable adversarial interventions. Per decision corpus review Decision 5 (HIGH confidence): under Option D+, the framework is a white-box interpretability probe (Gemini Part 5 "high-precision diagnostic probe", Part 4 "successful interpretability discovery"), not an adversarial-discovery tool. Rewrite preserves a brief dual-use acknowledgement (corpus mild caveat) + mirrors §7.1's already-applied white-box-access phrasing for stylistic consistency + frames future-research-program directionality. Bailey 2024 citation deferred per Josh sign-off ("defer the mild polish question"). — pending audit]
Reversible KV-Cache MCTS is a controlled-access white-box probe of the latent geometry of frozen language models; it requires white-box access and is not an inference-time attack vector against deployed systems. The Dimensional Escape pathology---that low-rank surrogate signals admit degenerate optima in high-dimensional cache space---informs the design of both monitoring procedures and test-time scaling schemes. We disclose findings openly; the open questions our results raise (which reward designs resist Dimensional Escape; whether monitor-driven Dimensional Escape manifests under inference-time KV-cache perturbation, Hypothesis~\ref{hyp:monitor_driven}) are inputs to robust monitor and reward-model design, not adversarial-discovery enablers.

\subsection{Limitations}

% [REVISION | §7.3-Limitations-recast-D+ | 2026-05-11 | Wholesale rewrite of §7.3 Limitations under Option D+ framing (Session A v9 draft 2026-05-10, applied 2026-05-11). Deletions from Option α+ version: (a) OE-hypothesis-not-yet-validated — superseded by Track F validation of Hypothesis 2 + sub-case structure under Definition def:dim_escape; (b) OEI-false-negatives — OEI is no longer §1 contribution under Option D+; (c) validation-chain-to-deceptive-intent — Option D+ explicitly decouples from intent, replaces with latent-response-under-steering. Rewrites: H-Neuron framing (latent-response diagnostic, not deception detector); σ_H framing (paired-with-entropy). Preserves: Procrustes uncertainty, RepE 1B-out-of-range, single-Llama-family. Adds 6 new Option D+ limitations: single model family, MCQ-letter-prompting under-statement (~20pp per Track C), single steering direction, 4-prompt-class Cartography scope, entropy-as-falsified-design framing, 200-cycle reversibility empirical bound. Net ~+410 words on §7.3 — Day-10 page-budget pass may reclaim by deferring items 5/6/8 to §7.5 Future Work. The 2026-04-24 `G1` marker is subsumed. — pending audit]
We acknowledge several limitations of the work as submitted. First, our empirical study covers only the Llama-3.2 model family at 1B and 3B parameter scales. While we observe Dimensional Escape via two structurally distinct routes at the two scales (1B prior-amplification, 3B greedy-collapse), we make no claim about MoE architectures (e.g., gpt-oss-20b), architectures outside the Llama family (e.g., Qwen, Mistral, Gemma), or scales beyond 3B. Generalization across model families is left to Phase B research; the algorithmic infrastructure of Reversible MCTS is architecture-agnostic by construction, but the Goodhart pathology we characterize may manifest with different specifics (different prior tokens, different collapse modes) in other models.

Second, our headline accuracy measurements (Track F, Experiment~1) use letter-prompted multiple-choice scoring on ARC-Easy. Greedy baseline accuracy with this protocol is $0.295$ at 1B and $0.845$ at 3B, both under-stated by approximately $20$ percentage points relative to published log-probability scoring of the same benchmark. Letter-prompted scoring is the natural protocol for our setting because entropy-MCTS steers the next-letter logit directly, so the comparison between MCTS arms is internally fair, but absolute baselines should not be compared cross-protocol. Log-probability scoring of MCTS-augmented outputs is a separate study deferred to Phase B.

% [REVISION | §7.3-Limitations-trim-Day10-batch3 | 2026-05-15 | Day-10 Batch 3 page-budget cut #10: trimmed §7.3 from 8 numbered items to 5. Items 1 (Llama-3.2 single family), 2 (letter-prompted MCQ caveat), 4 (σ_H/ρ_R diagnostic framing) preserved verbatim. Items 3 (single steering direction), 5 (4-prompt-class Cartography scope), 6 (200-cycle reversibility empirical bound) consolidated into one new methodological-constraints item. Item 7 (entropy-falsified-specific scoping) dropped: overlaps with Contribution 3 framing and §\ref{sec:exp1} Interpretation. Item 8 (Procrustes deferral + sub-7B probe validation) preserved as new item 5. Per Day-9 page-budget cut plan §3 rank #10 (MEDIUM risk, save 0.40 pp). — pending audit]
Third, methodological scope is constrained along three dimensions: (a)~the steering direction is a single random unit vector shared across items and MCTS arms (the COCONUT defense, §\ref{sec:exp1}); orthogonal, adversarial, and learned vectors are untested. (b)~Cartography (Experiment~2) covers four prompt classes ($n \le 100$ each); out-of-distribution prompts (adversarial probes, multilingual inputs, code, multi-turn CoT) are not characterized. (c)~Reversibility (Theorem~\ref{thm:reversibility}) is empirically validated over 200 apply-revert cycles; very-deep searches (depth $\ge 50$ or cycle counts $\gg 200$) are unvalidated empirically, though the theorem implies exact reversibility independent of depth.

% [REVISION | §7.3-Limitation4-delete-Day10-batch5 | 2026-05-15 | Day-10 Batch 5 drop-C2: DELETED former Limitation 4 (σ_H and ρ_R serve as latent-response diagnostics rather than as deception detectors). Under Batch 5 drop-C2 the telemetry-matrix framework is no longer a §1 paper-level contribution; it survives only as supplementary content in Appendix~\ref{sec:appendix-telemetry}. A limitation about "the deception-detection scope of our paper-level monitoring substrate" is therefore moot — the paper no longer presents that substrate as a contribution. The methodological caveat about H-Neuron extension to strategic deception is captured by the CRITIQUE NOTE in Appendix~\ref{sec:appendix-telemetry} (Bottom-Up Channel paragraph). Renumbers: former Limitation 5 (Procrustes / sub-7B probe validation) becomes new Limitation 4. — pending audit]
Fourth, cross-model steering-vector transfer via Procrustes alignment between model families (e.g., Llama 1B $\leftrightarrow$ Llama 7B, or Llama 3.2 $\leftrightarrow$ Mistral 7B) is an open research question with uncertain guarantees and is not addressed in this submission. The Representation Engineering probe validation of Zou et al.\ (2023)~\cite{zou2023repe} was conducted at $\ge 7$B parameter scales; our use of 1B and 3B scales for primary experiments falls below this validation range, and probe-direction reliability at sub-7B scale is an additional open empirical question that our results do not resolve.
\label{sec:limitations}

% [REVISION | §7.4-move-Day10-batch3 | 2026-05-15 | Day-10 Batch 3 page-budget cut #1 (largest single move): MOVED §7.4 Empirical Measurement Prerequisites and Initial Observations + table tab:oei-alpha-sweep to §E Measurement Pipeline in the appendix. The two-issue measurement-pipeline narrative (W_K projection issue + σ_H per-layer calibration) was honest hedging but read as bug-fix history rather than scientific content — camera-ready material best deferred to §E. Cross-refs to sec:measurement-prereqs and tab:oei-alpha-sweep resolve to the appendix after MOVE; the §5.5 OEI bullet reference was updated to "Appendix~\ref{sec:measurement-prereqs}" in this same batch. Five prior §7.4 layered REVISION markers (§7-measurement-prereqs 2026-05-06 AM, §7.4-post-fix-update 2026-05-06 PM, §7.4-σH-verification-update 2026-05-06 EOD, §7.4-σH-investigation-update 2026-05-07, §7.4-σH-resolved 2026-05-07; plus §7.4-opening-soften-D+ 2026-05-11 and OEI-alpha-sweep-relocate-to-§7.4-D+ 2026-05-12 and Exp1-table-σH-TDS-update 2026-05-07) are subsumed here. Per Day-9 page-budget cut plan §3 rank #1 (LOW risk, save 1.20 pp). — pending audit]
\subsection{Future Work}\label{sec:future-work}
% [REVISION | §7.5-FutureWork-trim-Day10-batch3 | 2026-05-15 | Day-10 Batch 3 page-budget cut #6: trimmed §7.5 from 4 paragraphs (Infrastructure / Monitor-design / Reward-design with 8-candidate enumeration / Cartography extensions) to 1 consolidated paragraph. The 8-candidate anti-Goodhart enumeration was the Pass K speculation-vs-page-budget tradeoff candidate; consolidated to "multi-channel ensembles, external PRMs, and self-consistency variants" representative shortlist. Monitor-design research paragraph reframed to align with Batch 5 drop-C2 (no telemetry-matrix-as-substrate framing). The 2026-05-11 §7.5-FutureWork-recast-D+ marker is subsumed; the venue-check CRITIQUE NOTE is dropped (ACL/EMNLP commitment locked). Per Day-9 page-budget cut plan §3 rank #6 (LOW risk, save 0.35 pp). — pending audit]

Phase~B directions: multi-GPU scaling for $M_{\text{KV}} \ge 20$~GB regimes ($\approx 162$~GB at 20B exceeds a single 80~GB H100); Procrustes-aligned steering-vector transfer; gpt-oss-20b MoE; sparse FP32 accumulators; anti-Goodhart reward designs surveyed in §\ref{sec:discussion-implications}; Cartography extensions to additional prompt classes and model families (§\ref{sec:exp2}). A Croissant~1.1 metadata package ships as supplementary material.

% [REVISION | §7.5-compress-Day14-Tianyu | 2026-05-18 | Compressed J2 license attribution + J3 AI assistant disclosure to two single-sentence statements. The "Built with Llama" prominent display + verbatim license notice text are Meta requirements for *model distribution*, not paper-text citation; one-sentence form with citation is standard ACL practice for AUP attribution. AI disclosure preserved per ARR Responsible NLP Checklist E1 (which only requires declaring AI use, not extended description).]
\subsection{Ethical Considerations}
This work uses Meta's Llama 3.2 under the Llama 3.2 Community License~\citep{meta2024llama3_2_aup}. AI assistants used in the execution window: Anthropic's Claude (coding/writing); Google's Gemini (one reward-design strategy review). All AI outputs were verified against primary artifacts; no empirical claims rest on AI-generated summaries.

% [REVISION | §8-cut | 2026-04-24 | Timeline and Milestones section removed entirely — research proposal artifact, not a paper section; ~0.5 page saved — pending audit]
